\documentclass[11pt,a4paper]{article}
\usepackage[utf8]{inputenc}
\usepackage[T1]{fontenc}
\usepackage[spanish]{babel}
\usepackage{hyperref}
\usepackage{geometry}
\geometry{margin=2.5cm}
\usepackage{parskip}

\title{349 esto no es para la academia, es para el mundo}
\author{Alejandro Valenzuela Rivera\\\small Universidad Autónoma de Nuevo León}
\date{\today}
\begin{document}
\maketitle

\begin{quote}
\textbf{Sesgo declarado.} Este trabajo se sostiene desde el sesgo declarado de su autor, y no finge neutralidad. Lo que afirma se puede tumbar, y ese es el punto: el falsificador está escrito, no escondido.
\end{quote}

\section*{Resumen}
Esto no se escribe para la academia: se escribe para el mundo. Y el mundo tiene la academia quieras o no, así que la academia lo tendrá que leer.

\section{Por qué se sostiene}
Lo que se escribe para un gremio se lee en el gremio; lo que se escribe para el mundo obliga al gremio a leerlo.

\section{De dónde viene}
La autoridad: 'no es para la academia, es para el mundo, y el mundo tiene la academia quieras o no'.

\section{El caso}
El repo: lo va a leer la academia, y no se escribió para ella.

\section{Método}
\textquestiondown{}El que no es de la academia lo puede usar sin pedir permiso?

\section{Límites, y qué lo tumba}
Una obra que solo se entiende con el vocabulario de un gremio.

\section{Lo que queda abierto}
\textquestiondown{}Se escribió para que lo lea el gremio, o para que lo lea el mundo?

\section*{Cómo se leyó esta tesis}
Esto se escribió en el repositorio man-made-empathy, del que este candado es una pieza: candados/04 el lenguaje/349 esto no es para la academia, es para el mundo.md. La máquina que lo compuso lo hizo desde este mismo archivo, y no agregó nada que no estuviera escrito. Si algo de esto es falso, se dice en el mismo sitio, y se corrige con fecha.

\end{document}
